\documentclass[11pt]{article}
\usepackage[margin=1in]{geometry}
\usepackage[T1]{fontenc}
\usepackage[utf8]{inputenc}
\usepackage{lmodern}
\usepackage{amsmath,amssymb,amsthm}
\usepackage{array}
\usepackage{booktabs}
\usepackage{longtable}
\usepackage{enumitem}
\usepackage{hyperref}
\usepackage{xcolor}

\hypersetup{
  colorlinks=true,
  linkcolor=blue,
  urlcolor=blue,
  citecolor=blue
}

\setlist[itemize]{topsep=2pt,itemsep=2pt,parsep=0pt,leftmargin=1.5em}
\setlist[enumerate]{topsep=2pt,itemsep=2pt,parsep=0pt,leftmargin=1.7em}

\newcommand{\code}[1]{\texttt{\detokenize{#1}}}
\newcommand{\paperref}[1]{\textbf{Paper ref:} \code{#1}}
\newcommand{\coderef}[1]{\textbf{Code ref:} \code{#1}}
\newcommand{\status}[1]{\textbf{Status:} #1}
\newcommand{\explicit}{\textbf{Explicitly stated.}}
\newcommand{\implicit}{\textbf{Implicitly used.}}
\newcommand{\inferred}{\textbf{Inferred.}}
\newcommand{\missing}{\textbf{Missing.}}

\title{Technical Diagnosis of the L2D-SLDS Paper and Implementation}
\author{Repository audit written from \code{Paper/icml2026/main.tex} and local code}
\date{2026-03-25}

\begin{document}
\maketitle

\begin{abstract}
This document is a rigorous diagnosis of the research paper in
\code{Paper/icml2026/main.tex} and the accompanying implementation in the
repository root. The goal is not to summarize at a high level, but to
reconstruct the exact mathematical object being proposed, identify what is and
is not formally defined, map the implementation back to the paper line by line,
and isolate paper errors, code errors, hidden assumptions, and protocol
mismatches. Every important criticism below is justified by direct references to
the paper or code.
\end{abstract}

\tableofcontents

\section{Paper Overview}

\subsection{Problem setting}

The paper studies sequential expert routing under partial feedback,
non-stationarity, and dynamic expert availability. The formal setup is given in
\code{Paper/icml2026/Section/Approach.tex:10--106}. At round $t$:
\begin{itemize}
    \item the environment reveals a context $x_t \in \mathcal{X} \subseteq \mathbb{R}^d$,
    \item a target $y_t \in \mathcal{Y} \subseteq \mathbb{R}^{d_y}$ exists,
    \item a non-empty available expert set $\mathcal{E}_t$ is observed,
    \item the router selects $I_t \in \mathcal{E}_t$,
    \item only the queried prediction $\hat y_{t,I_t}$ and target $y_t$ are observed,
    \item therefore only the queried residual
    \[
    e_t = e_{t,I_t} = \hat y_{t,I_t} - y_t
    \]
    and queried cost $C_{t,I_t}$ are observed.
\end{itemize}

The paper defines the potential residual and cost of expert $k$ at time $t$ as
\[
e_{t,k} = \hat y_{t,k} - y_t, \qquad
C_{t,k} = \psi(e_{t,k}) + \beta_k,
\]
where $\psi$ is a convex loss and $\beta_k \ge 0$ is a known consultation fee;
see \code{Paper/icml2026/Section/Approach.tex:24--43}.

The paper allows a time-varying expert registry $\mathcal{K}_t$ that may be
larger than the currently available set $\mathcal{E}_t$, and stale experts may
be pruned from the registry; see
\code{Paper/icml2026/Section/Approach.tex:17--20} and
\code{Paper/icml2026/Section/Approach.tex:343--447}.

\subsection{Main claim}

The main claim is that a factorized switching linear dynamical system over
expert residuals yields a tractable predictive belief over \emph{unqueried}
experts under censoring, and that this can be used for routing by minimizing a
myopic IDS-style information ratio. This claim is developed in
\code{Paper/icml2026/Section/Approach.tex:108--340}. The latent structure has:
\begin{itemize}
    \item a regime variable $z_t$,
    \item a shared factor $g_t$,
    \item expert-specific states $u_{t,k}$,
    \item context-dependent regime transitions $\Pi_\theta(x_t)$,
    \item a residual emission model
    \[
    e_{t,k} \mid z_t, g_t, u_{t,k}, x_t \sim
    \mathcal{N}\!\left(\Phi(x_t)^\top(B_k g_t + u_{t,k}), R_{z_t,k}\right).
    \]
\end{itemize}

The claimed benefit of the shared factor is cross-expert information transfer:
after querying one expert, posterior beliefs about other experts change because
their residual laws share dependence on $g_t$.

\subsection{Claimed contributions}

The claimed contributions are listed in
\code{Paper/icml2026/Section/Introduction.tex:28--36}:
\begin{enumerate}
    \item a formulation of online expert routing for non-stationary time series
    with bandit feedback and dynamic availability;
    \item the L2D-SLDS model;
    \item an IMM-style filtering recursion and pruning mechanism;
    \item an IDS-inspired routing rule based on information about $(z_t, g_t)$;
    \item empirical evaluation on synthetic, Melbourne, and FRED tasks.
\end{enumerate}

\paragraph{Assessment of the contribution statement.}
This list is internally coherent. However, it mixes three different levels of
claim:
\begin{itemize}
    \item formal problem definition,
    \item algorithm design,
    \item empirical evidence.
\end{itemize}
What it does \emph{not} provide is a learning-theoretic guarantee such as regret
or consistency. The framing invokes online learning and bandits, but the paper's
formal guarantee layer is much weaker than the framing suggests.

\section{Precise Reconstruction of the Method}

\subsection{Formal setup}

I separate here what the paper states explicitly from what is only implied.

\explicit The decision-time sigma-algebra is
\[
\mathcal{F}_t = \sigma(\mathcal{H}_{t-1}, x_t, \mathcal{E}_t),
\]
with $\mathcal{H}_t = ((x_\tau,\mathcal{E}_\tau,O_\tau))_{\tau=1}^t$ and
$O_t = (I_t, \hat y_{t,I_t}, y_t)$; see
\code{Paper/icml2026/Section/Approach.tex:52--63}.

\explicit A policy $\pi_t(\cdot \mid \mathcal{F}_t)$ is an
$\mathcal{F}_t$-measurable distribution over $\mathcal{E}_t$, and the action is
sampled as $I_t \sim \pi_t(\cdot \mid \mathcal{F}_t)$; see
\code{Paper/icml2026/Section/Approach.tex:59--63}.

\explicit The objective is expected cumulative cost
\[
J(\pi) = \mathbb{E}\left[\sum_{t=1}^T C_{t,I_t}\right];
\]
see \code{Paper/icml2026/Section/Approach.tex:90--106}.

\explicit As a one-step ideal benchmark, the paper defines
\[
k_t^\star \in \arg\min_{k \in \mathcal{E}_t} \mathbb{E}[C_{t,k} \mid \mathcal{F}_t],
\]
which it calls the myopic Bayes selector; see
\code{Paper/icml2026/Section/Approach.tex:96--105}.

\missing The paper does \emph{not} define:
\begin{itemize}
    \item a regret notion,
    \item a comparator class over policies,
    \item a non-stationary regret benchmark,
    \item a bandit sample complexity target,
    \item or a consistency criterion for the router.
\end{itemize}

Therefore the problem is a sequential decision problem, but not a formal online
learning theorem statement.

\subsection{Variables and notation}

The core notation in the model section is mostly coherent:
\begin{itemize}
    \item $z_t \in \{1,\dots,M\}$: latent regime;
    \item $g_t \in \mathbb{R}^{d_g}$: shared latent factor;
    \item $u_{t,k} \in \mathbb{R}^{d_\alpha}$: expert-specific latent state;
    \item $B_k \in \mathbb{R}^{d_\alpha \times d_g}$: loading matrix for expert $k$;
    \item $\Phi: \mathcal{X} \to \mathbb{R}^{d_\alpha \times d_y}$: feature map;
    \item $\alpha_{t,k} = B_k g_t + u_{t,k}$: latent reliability vector;
    \item $R_{m,k} \in \mathbb{S}^{d_y}_{++}$: emission covariance.
\end{itemize}

\paragraph{Well-definedness.}
The expression
\[
\Phi(x_t)^\top \alpha_{t,k}
\]
is dimensionally valid because $\Phi(x_t)^\top \in \mathbb{R}^{d_y \times d_\alpha}$
and $\alpha_{t,k} \in \mathbb{R}^{d_\alpha}$, so the output lies in
$\mathbb{R}^{d_y}$; see
\code{Paper/icml2026/Section/Approach.tex:211--230}.

\paragraph{Notation issue.}
The paper alternates between $d_\alpha$ and the implementation's $d_\phi$ for
the idiosyncratic-state dimension. This is not mathematically wrong, but it
creates friction when mapping paper to code. The implementation uses
\code{d_phi} as the dimension of $u_{t,k}$ and of the feature vector returned by
\code{feature_phi}; see \code{models/factorized_slds.py:68--110} and
\code{models/router_model.py:5--9}.

\subsection{Objective}

The paper's operational decision rule is myopic. It computes predicted
one-step costs, then combines them with an exploration score. The exploitation
part is
\[
\bar C_{t,k}^{\mathrm{pred}}
=
\mathbb{E}\bigl[\psi(e_{t,k}^{\mathrm{pred}})\mid \mathcal{F}_t\bigr] + \beta_k
\]
and the predictive gap is
\[
\Delta_t(k)
=
\bar C_{t,k}^{\mathrm{pred}} - \bar C_{t,k_t^{\mathrm{pred}}}^{\mathrm{pred}},
\qquad
k_t^{\mathrm{pred}} \in \arg\min_k \bar C_{t,k}^{\mathrm{pred}}.
\]
See \code{Paper/icml2026/Section/Approach.tex:306--318} and
\code{Paper/icml2026/Section/AppendixParts/info_gain.tex:97--161}.

\paragraph{Important limitation.}
This is not a dynamic programming objective. It is a one-step predictive
criterion plus an information heuristic. The paper is explicit about this point
and does not claim Bayes optimal control; see
\code{Paper/icml2026/Section/Approach.tex:299--304}.

\subsection{Model}

\subsubsection{Regime dynamics}

The regime dynamics are context-dependent:
\[
\mathbb{P}(z_t = m \mid z_{t-1} = \ell, x_t)
=
\Pi_\theta(x_t)_{\ell m}.
\]
The paper parameterizes $\Pi_\theta(x_t)$ by low-rank attention logits:
\[
S(x_t) = \frac{1}{\sqrt{d_{\mathrm{attn}}}} Q_\theta(x_t) K_\theta(x_t)^\top,
\qquad
\Pi_\theta(x_t)_{\ell m}
=
\frac{\exp(S_{\ell m}(x_t))}{\sum_j \exp(S_{\ell j}(x_t))}.
\]
See \code{Paper/icml2026/Section/Approach.tex:135--163}.

\paragraph{Interpretation.}
\explicit The paper intends a learned context-dependent transition kernel.

\implicit This assumes enough data and a stable enough training problem to
identify the dependence of regime switching on context.

\missing The paper never states what happens when transition parameters are
untrained or weakly trained, even though the code permits exactly that.

\subsubsection{Continuous latent dynamics}

Conditioned on regime $z_t = m$:
\[
g_t = A_m^{(g)} g_{t-1} + w_t^{(g)}, \qquad
w_t^{(g)} \sim \mathcal{N}(0,Q_m^{(g)}),
\]
\[
u_{t,k} = A_m^{(u)} u_{t-1,k} + w_{t,k}^{(u)}, \qquad
w_{t,k}^{(u)} \sim \mathcal{N}(0,Q_m^{(u)}).
\]
See \code{Paper/icml2026/Section/Approach.tex:165--203}.

\paragraph{Comment.}
The paper shares $(A_m^{(u)},Q_m^{(u)})$ across experts and uses expert-specific
heterogeneity only through $u_{t,k}$ realizations and $B_k$; this is a modeling
choice that improves statistical strength but imposes a strong inductive bias.

\subsubsection{Emission model}

The latent reliability of expert $k$ is
\[
\alpha_{t,k} = B_k g_t + u_{t,k},
\]
and the residual emission is
\[
e_{t,k} \mid (z_t=m, g_t, u_{t,k}, x_t)
\sim
\mathcal{N}\!\left(\Phi(x_t)^\top \alpha_{t,k}, R_{m,k}\right).
\]
See \code{Paper/icml2026/Section/Approach.tex:205--239}.

\paragraph{Interpretation.}
\explicit The model is defined for residuals, not losses.

\inferred This is because residuals can be negative and therefore admit a linear
Gaussian state-space treatment, whereas losses generally do not.

\missing The paper does not discuss misspecification when the true residual law
is heavy-tailed, skewed, heteroskedastic in a way not captured by $R_{m,k}$, or
nonlinear in $\Phi(x_t)$.

\subsection{Inference / estimation}

\subsubsection{Filtering}

The filter is described in algorithmic form in
\code{Paper/icml2026/Section/AppendixParts/alg_router.tex:10--58}. The intended
per-round recursion is:
\begin{enumerate}
    \item update the registry $\mathcal{K}_t$ by adding current experts and
    pruning stale unavailable experts;
    \item initialize entering experts with default idiosyncratic priors;
    \item perform an IMM predictive step using $\Pi_\theta(x_t)$;
    \item time-update the Gaussian beliefs of $g_t$ and each stored $u_{t,k}$;
    \item compute per-expert predictive residual laws;
    \item choose an action via predicted cost and information ratio;
    \item after observing the queried residual, update the joint state
    $(g_t,u_{t,I_t})$ by a Kalman correction;
    \item project the updated covariance back to a factorized family by dropping
    posterior cross-covariances.
\end{enumerate}

\paragraph{Mathematical status.}
The exact queried update is standard Kalman conditioning on the joint state. The
approximation enters only at the projection step.

\subsubsection{Cross-covariance and factorization}

The appendix explicitly notes that even if the predictive covariance is
block-diagonal, the posterior covariance after conditioning on the queried
residual generally has a non-zero $g$--$u_{I_t}$ cross-covariance; see
\code{Paper/icml2026/Section/AppendixParts/derivations.tex:4--66}. This is
important:
\begin{itemize}
    \item the paper is aware that the exact posterior is not factorized;
    \item the factorized filter is therefore a projection approximation, not an
    exact Bayes update.
\end{itemize}

\paragraph{Critique.}
This is the central approximation of the method, but the paper gives no
quantitative control on the error induced by repeatedly discarding these
cross-covariances. That is acceptable for an empirical methods paper, but not
for a learning-theoretic claim.

\subsubsection{Parameter learning}

The paper's learning routine is optional MCEM; see
\code{Paper/icml2026/Section/AppendixParts/alg_learning.tex:9--112}. The
algorithmic idea is:
\begin{itemize}
    \item sample regime paths $z_{1:T}$ given current latent trajectories,
    \item sample $g_{1:T}$ given $z_{1:T}$ and expert-specific trajectories,
    \item sample $u_{1:T,k}$ expertwise,
    \item update $(A,Q,B,R)$ by weighted moment matching / regression,
    \item update the transition model parameters $\theta$ from expected regime
    transition counts.
\end{itemize}

\paragraph{Assessment.}
The algorithm is plausible and the code largely implements it. What is missing
is a statement of convergence assumptions or identifiability.

\subsection{Decision rule}

The decision rule is an IDS-inspired information ratio:
\[
I_t \in \arg\min_{k \in \mathcal{E}_t} \frac{\Delta_t(k)^2}{\mathrm{IG}_t(k)},
\]
where
\[
\mathrm{IG}_t(k)
=
\mathcal{I}\!\left((z_t,g_t); e_{t,k}^{\mathrm{pred}} \mid \mathcal{F}_t\right).
\]
The decomposition
\[
\mathcal{I}((z_t,g_t);e)
=
\mathcal{I}(z_t;e)
+
\mathcal{I}(g_t;e \mid z_t)
\]
is given in \code{Paper/icml2026/Section/AppendixParts/info_gain.tex:7--34} and
\code{Paper/icml2026/Section/AppendixParts/info_gain.tex:163--260}.

\paragraph{Shared-factor term.}
The closed-form term
\[
\frac{1}{2}\log\det\!\left(I + H_{t,k}\Sigma_{g,t|t-1}^{(m)} H_{t,k}^\top
(S_{t,k}^{(m)})^{-1}\right)
\]
is correct for a linear Gaussian channel.

\paragraph{Mode-identification term.}
The paper estimates $\mathcal{I}(z_t;e_{t,k}^{\mathrm{pred}} \mid \mathcal{F}_t)$
by Monte Carlo under a Gaussian mixture model. This is technically sound as an
estimator of the quantity it defines.

\paragraph{Theoretical status.}
\explicit The paper states this is a pragmatic exploration heuristic, not an
optimality guarantee; see
\code{Paper/icml2026/Section/Approach.tex:299--304}.

\paragraph{Critical point.}
This honesty matters. The paper should not be read as proving an IDS guarantee.
It does not. It only borrows the information-ratio shape.

\subsection{Computational complexity}

The paper mentions scalability motivations, but it does not provide a clean
complexity statement. The practical complexity depends on:
\begin{itemize}
    \item the size of the active registry,
    \item the latent dimensions $d_g$ and $d_\alpha$,
    \item the number of regimes $M$,
    \item the Monte Carlo budget for $I(z_t;e)$.
\end{itemize}

\paragraph{What is clear.}
The factorized representation avoids storing a full joint covariance across all
experts, which would scale poorly with registry size.

\paragraph{What is missing.}
There is no explicit asymptotic runtime or memory claim in the paper, so the
reader cannot evaluate the scalability gain precisely from the manuscript alone.

\section{Assumptions and Hidden Assumptions}

\subsection{Explicit assumptions}

The following assumptions are stated or directly encoded in the formal model:
\begin{itemize}
    \item \textbf{Exogeneity of the data stream:} past routing actions do not
    affect the law of $(x_t,\mathcal{E}_t,y_t)$ beyond the censoring of
    observations; see \code{Paper/icml2026/Section/Approach.tex:76--89}.
    \item \textbf{Conditional Gaussianity:} both latent dynamics and emissions
    are linear Gaussian given the regime; see
    \code{Paper/icml2026/Section/Approach.tex:174--239}.
    \item \textbf{Positive-definite noise covariances:} $Q_m^{(g)}$,
    $Q_m^{(u)}$, and $R_{m,k}$ are positive definite.
    \item \textbf{Factorized filtering approximation:} posterior
    cross-covariances are discarded after each queried update; see
    \code{Paper/icml2026/Section/Approach.tex:250--257} and
    \code{Paper/icml2026/Section/AppendixParts/derivations.tex:57--66}.
    \item \textbf{Shared idiosyncratic dynamics parameters across experts:}
    $(A_m^{(u)},Q_m^{(u)})$ do not depend on $k$.
\end{itemize}

\subsection{Implicit assumptions}

The method also relies on several assumptions that are not made explicit enough:
\begin{itemize}
    \item \textbf{Persistent expert identities.} The meaning of re-entry depends
    on an expert ID corresponding to the same forecasting mechanism over time.
    This is used operationally but not formally discussed.
    \item \textbf{Identifiability of the decomposition} $B_k g_t + u_{t,k}$.
    Without additional constraints, there can be tradeoffs between the shared and
    idiosyncratic components.
    \item \textbf{Stable enough dynamics.} The filter and EM updates become
    difficult to interpret if $A_m^{(g)}$ or $A_m^{(u)}$ imply explosive latent
    trajectories.
    \item \textbf{Availability process does not induce confounding.} If
    availability is itself coupled to latent difficulty in a way that violates
    exogeneity, the probabilistic interpretation changes.
    \item \textbf{Feature map adequacy.} The model assumes $\Phi(x_t)$ captures
    the context dependence of residuals well enough for linear-Gaussian modeling
    to be informative.
\end{itemize}

\subsection{Missing assumptions}

Several assumptions are needed for stronger claims but are absent:
\begin{itemize}
    \item conditions for identifiability of $(B_k,g_t,u_{t,k})$;
    \item consistency or convergence conditions for MCEM under censoring;
    \item approximation-error control for repeated posterior projection;
    \item assumptions under which the IDS-inspired rule improves long-run
    performance;
    \item assumptions under which pruning and re-entry preserve not just retained
    marginals, but decision quality for re-entering experts.
\end{itemize}

\section{Theory Audit}

\subsection{Evaluation of the formal claims}

\subsubsection{Problem formulation}

The problem formulation is precise enough to define the stochastic process and
the one-step Bayes benchmark. However, it does not support a specialist-level
online learning claim in the regret sense, because no comparator or regret is
defined.

\paragraph{Diagnosis.}
The paper is better described as a model-based sequential decision method under
partial feedback than as a learning-theoretic online-learning paper.

\subsubsection{Proposition on cross-update}

Proposition \code{prop:cross_update} states that, conditional on fixed regime and
joint Gaussianity, observing the queried residual changes the predictive mean of
another expert if and only if the predictive cross-covariance is non-zero; see
\code{Paper/icml2026/Section/Approach.tex:259--278} and
\code{Paper/icml2026/Section/AppendixParts/proofs.tex:10--38}.

\paragraph{Assessment.}
This proposition is correct. It is a direct consequence of Gaussian
conditioning:
\[
\mathbb{E}[X \mid Y=y]
=
\mu_X + \Sigma_{XY}\Sigma_Y^{-1}(y-\mu_Y).
\]

\paragraph{Limitation.}
It is an elementary structural fact, not a performance theorem.

\subsubsection{Pruning proposition}

Proposition \code{prop:invariance} states that pruning experts from the registry
does not change the retained marginals because pruning is defined by
marginalization; see
\code{Paper/icml2026/Section/Approach.tex:384--405} and
\code{Paper/icml2026/Section/AppendixParts/proofs.tex:44--86}.

\paragraph{Assessment.}
The proposition is correct but tautological. It proves exactly what its
definition encodes.

\paragraph{What it does not prove.}
It does not show that re-entered experts are well initialized, only that
retained experts are not perturbed by pruning.

\subsubsection{Birth coupling proposition}

Proposition \code{prop:transfer} states that under the factorized predictive
belief,
\[
\mathrm{Cov}(\alpha_{t,j}, \alpha_{t,k} \mid \mathcal{F}_t,z_t=m)
=
B_j \Sigma_{g,t|t-1}^{(m)} B_k^\top,
\]
so shared-factor uncertainty induces coupling across experts; see
\code{Paper/icml2026/Section/Approach.tex:449--464} and
\code{Paper/icml2026/Section/AppendixParts/proofs.tex:92--117}.

\paragraph{Assessment.}
This proposition is correct as stated for the latent reliability vectors
$\alpha_{t,j},\alpha_{t,k}$.

\paragraph{Important caveat.}
It is \emph{not} the proposition that justifies queried update transfer at the
residual level. That role belongs to \code{prop:cross_update}. The paper later
mis-cites this point in the synthetic results.

\subsubsection{IDS-inspired exploration}

The paper's mutual-information construction is technically coherent and the
shared-factor term is derived correctly. The mode-identification term is
estimated by Monte Carlo because the predictive residual law is a Gaussian
mixture; see
\code{Paper/icml2026/Section/AppendixParts/info_gain.tex:31--34} and
\code{Paper/icml2026/Section/AppendixParts/info_gain.tex:204--260}.

\paragraph{Assessment.}
As a heuristic exploration score under the model, this is fine.

\paragraph{What is not justified.}
There is no proof that minimizing the resulting ratio improves cumulative cost,
minimizes Bayes regret, or satisfies any IDS theorem. The paper is careful in
its wording, but the framing is still stronger than the theorem layer.

\subsection{Proof gaps or questionable steps}

\begin{enumerate}
    \item \textbf{No approximation analysis for factorization.}
    The main computational approximation is the repeated projection to a
    factorized posterior. The paper acknowledges this, but gives no bound on how
    much predictive distributions or information scores may drift from the exact
    posterior.

    \item \textbf{Latent-to-residual transfer is not fully unpacked.}
    The paper correctly notes that latent shared-factor covariance induces
    residual cross-covariance through $\Phi(x_t)^\top B_k$, but it does not spell
    out the minimal non-degeneracy conditions under which this actually affects
    the queried conditional mean for another expert.

    \item \textbf{No identifiability discussion.}
    The same residual variation can often be explained by $g_t$ or by
    $u_{t,k}$. This matters both conceptually and for MCEM, but the paper does
    not discuss it.
\end{enumerate}

\subsection{Claims correct only under stronger assumptions}

\begin{itemize}
    \item \textbf{Context-aware switching.} This only means something if
    transition parameters are trained or otherwise intentionally specified. In
    the actual synthetic and Melbourne configs, EM is disabled while the
    transition mode is attention-based, so the model uses random transition
    parameters unless trained elsewhere.

    \item \textbf{General vector output $d_y \ge 1$.} The notation supports it,
    but all provided environments and experiments are scalar-output. The claim is
    formal generality, not demonstrated generality.

    \item \textbf{Cross-expert transfer in practice.} The theoretical mechanism
    requires non-trivial shared loadings and a feature map that does not suppress
    the shared channel. Without these, transfer can vanish.
\end{itemize}

\subsection{Possible counterexamples or failure modes}

\begin{itemize}
    \item If $B_k = 0$ for all $k$, then the shared factor plays no role and the
    model collapses to independent expert-specific processes.
    \item If $\Phi(x_t)$ is zero or orthogonal to the shared component, then even
    non-zero latent shared covariance need not produce residual-level transfer.
    \item If availability depends adversarially on latent difficulty, the
    exogeneity assumption fails and the interpretation of partial feedback as
    pure censoring becomes invalid.
    \item If the true residual law is strongly non-Gaussian, the information
    scores may be badly calibrated.
\end{itemize}

\section{Implementation Walkthrough}

\subsection{Repository structure}

The implementation relevant to the paper is distributed across:
\begin{itemize}
    \item \code{slds_imm_router.py}: experiment driver and router construction;
    \item \code{models/factorized_slds.py}: proposed model implementation;
    \item \code{router_eval.py}: interaction loops with environments;
    \item \code{plot_utils.py}: aggregation, masking, and reporting;
    \item \code{environment/synthetic_env.py}: synthetic benchmark;
    \item \code{environment/etth1_env.py}: time-series data environment also used
    for FRED-like and Melbourne-like experiments;
    \item \code{models/linucb_baseline.py} and
    \code{models/neuralucb_baseline.py}: baseline implementations;
    \item \code{config/exp_synthetic_1.yaml},
    \code{config/config_melbourne.yaml}, and \code{config/exp_FRED.yaml}:
    experiment protocols.
\end{itemize}

\subsection{Entry points}

The proposed method is instantiated in
\code{slds_imm_router.py:1247--1575}. The driver builds one or more
\code{FactorizedSLDS} instances depending on:
\begin{itemize}
    \item feedback mode,
    \item transition mode,
    \item exploration mode,
    \item inclusion of the no-$g_t$ ablation.
\end{itemize}

Offline EM, when enabled, is dispatched from
\code{slds_imm_router.py:1980--2104} and parameter copies are shared across
paired router instances in \code{slds_imm_router.py:2295--2367}.

\subsection{Data flow}

At evaluation time, \code{router_eval.py} does the following:
\begin{enumerate}
    \item fetch context and availability at time $t$,
    \item call the router's \code{select\_expert},
    \item obtain either residual or loss observations,
    \item record cost,
    \item update the router.
\end{enumerate}
See \code{router_eval.py:486--554}.

For \code{FactorizedSLDS}, the observation passed to the model is the residual,
not the loss, in \code{router_eval.py:216--243}. This matches the paper.

\subsection{Model components}

\subsubsection{Observation model in code}

The class docstring in \code{models/factorized_slds.py:48--66} implements the
same residual emission model as the paper:
\[
e_{t,k} = \phi(x_t)^\top(B_k g_t + u_{t,k}) + \varepsilon.
\]

The implementation only supports residual observations:
\code{models/factorized_slds.py:107--108} and
\code{models/factorized_slds.py:259--264}.

\subsubsection{Feature map}

The default feature map is identity:
\code{models/router_model.py:5--9}.

\paragraph{Consequence.}
For the provided experiments, $\Phi(x_t)$ is effectively a vectorized identity
map on the context features. This is simpler than the general matrix-valued
$\Phi$ notation in the paper.

\subsubsection{Transition model}

The implementation supports:
\begin{itemize}
    \item an attention-based transition parameterization,
    \item a linear transition parameterization,
    \item an optional neural transition model.
\end{itemize}
The default is
\code{transition_mode="attention"} and \code{transition_init="random"} in
\code{models/factorized_slds.py:114--115}. The decision-time transition matrix
is produced by \code{\_context\_transition} in
\code{models/factorized_slds.py:741--763}.

\paragraph{Critical implementation fact.}
If EM is disabled and no trained transition model is loaded, then
\code{\_context\_transition} uses random initialized parameters. This is not a
minor implementation detail; it affects the meaning of the paper's
``context-aware switching'' claim in the no-EM experiments.

\subsection{Training loop}

The EM implementation in \code{models/factorized_slds.py:3390--4050} broadly
matches the appendix:
\begin{itemize}
    \item latent regime sampling,
    \item Kalman sampling for $g_t$,
    \item Kalman sampling for each $u_{t,k}$,
    \item M-step updates for $A_g$, $Q_g$, $A_u$, $Q_u$, $B_k$, $R_{m,k}$,
    \item gradient-based transition training.
\end{itemize}

\paragraph{Positive note.}
This is one of the stronger paper-to-code alignments in the repository.

\subsection{Inference / prediction}

Predictive statistics are computed in
\code{models/factorized_slds.py:1377--1459}. The code:
\begin{itemize}
    \item computes per-mode predictive means and covariances,
    \item computes expected loss under each mode,
    \item computes the shared-factor mutual-information term,
    \item returns per-expert predicted cost and information gain.
\end{itemize}

Mode-identification information is estimated in
\code{models/factorized_slds.py:1535--1613}. IDS selection is implemented in
\code{models/factorized_slds.py:1679--1786}.

\subsection{Evaluation procedure}

Evaluation masking occurs in \code{plot_utils.py:1085--1155}: costs and choices
before \code{em\_tk} are replaced by \code{NaN}. Reported average costs are then
computed with \code{np.nanmean} in \code{plot_utils.py:1291--1335}. Constant
expert baselines are also recomputed starting from \code{t\_start = em\_tk + 1}
in \code{plot_utils.py:1267--1290}.

\paragraph{Interpretation.}
This means the reported summaries for experiments with EM warm-up are
\emph{post-warmup} averages, not full-horizon averages, unless explicitly stated
otherwise.

\section{Paper-to-Code Alignment}

\small
\begin{longtable}{p{0.18\textwidth} p{0.19\textwidth} p{0.20\textwidth} p{0.12\textwidth} p{0.24\textwidth}}
\toprule
\textbf{Paper component} & \textbf{Location in paper} & \textbf{Code file/function} & \textbf{Status} & \textbf{Explanation} \\
\midrule
\endfirsthead
\toprule
\textbf{Paper component} & \textbf{Location in paper} & \textbf{Code file/function} & \textbf{Status} & \textbf{Explanation} \\
\midrule
\endhead

Sequential partial-feedback routing
& \code{Approach.tex:45--73}
& \code{router_eval.py:216--243}, \code{router_eval.py:486--554}
& Match
& The code passes only the queried residual to \code{FactorizedSLDS} in partial mode, matching the paper's observation model. \\

Residual state-space model
& \code{Approach.tex:205--239}
& \code{models/factorized_slds.py:48--66}
& Match
& The implemented observation model is residual-based and uses the same shared plus idiosyncratic decomposition. \\

Context-dependent regime transitions
& \code{Approach.tex:135--163}
& \code{models/factorized_slds.py:741--763}
& Partial match
& The code implements context-dependent transitions, but in some published experiments the parameters are random because EM is disabled. \\

Predictive cost computation
& \code{info_gain.tex:128--161}
& \code{models/factorized_slds.py:1377--1459}
& Match
& The code computes per-mode predictive moments and expected loss exactly in the spirit of the appendix. \\

Shared-factor IG term
& \code{info_gain.tex:188--202}
& \code{models/factorized_slds.py:1431--1443}
& Match
& The scalar and multivariate log-det forms are implemented. \\

Mode-identification term
& \code{info_gain.tex:204--260}
& \code{models/factorized_slds.py:1535--1613}
& Match
& Monte Carlo KL-to-mixture estimator is implemented with stable log-sum-exp. \\

IDS-inspired routing rule
& \code{Approach.tex:333--340}
& \code{models/factorized_slds.py:1679--1786}
& Match with extra heuristics
& The code adds score clipping and greedy fallbacks not documented in the paper. \\

Queried Kalman update and factorized projection
& \code{alg_router.tex:39--58}
& \code{models/factorized_slds.py:1980--2068}
& Match
& The code updates the joint state $(g_t,u_{t,I_t})$, then stores only block marginals. \\

Registry pruning
& \code{Approach.tex:362--405}
& \code{models/factorized_slds.py:2385--2413}
& Match
& Stale unavailable experts are dropped after \code{Delta_max}; re-entry restarts from default priors. \\

Offline MCEM
& \code{alg_learning.tex:9--112}
& \code{models/factorized_slds.py:3390--4050}
& Match
& The code contains a genuine latent-variable EM-style procedure broadly consistent with the appendix. \\

Synthetic experiment design
& \code{Experiments.tex:55--123}
& \code{environment/synthetic_env.py:547--582}, \code{config/exp_synthetic_1.yaml:44--79}
& Mismatch
& The code uses a latent model-matched residual generator with \code{shared_dim: 2}, not the additive expert-noise simulator with \mbox{$d_g=1$} described in the paper. \\

Synthetic warmup of 100 steps
& \code{Experiments.tex:121--123}
& \code{router_eval.py:514--554}, \code{config/exp_synthetic_1.yaml:98--100}
& Mismatch
& The published synthetic config disables EM and the interaction starts immediately from $t=1$. \\

Melbourne no-EM claim
& \code{Experiments.tex:158--162}
& \code{config/config_melbourne.yaml:123--153}
& Match at config level
& EM is indeed disabled in the Melbourne config. However, the transition model remains attention-based and therefore untrained unless separately learned. \\

FRED evaluation protocol
& Appendix experiment text
& \code{config/exp_FRED.yaml:166--168}, \code{plot_utils.py:1085--1155}
& Mismatch
& The code uses offline full-feedback EM on the prefix and reports post-prefix averages after masking the warmup. \\

NeuralUCB baseline
& \code{AppendixParts/Experiments.tex:83--110}
& \code{models/neuralucb_baseline.py:5--25}
& Mismatch
& The paper describes gradient-feature NeuralUCB; the code explicitly implements a NeuralLinear-style approximation. \\

General $d_y \ge 1$ support
& \code{Approach.tex:13--15}, \code{Approach.tex:211--230}
& \code{models/router_model.py:5--9}, environments
& Unclear / weakly supported
& The code can accept vector residuals in principle, but all provided experiments are scalar-output and the default feature map is vector-valued identity. \\

\bottomrule
\end{longtable}
\normalsize

\section{Issues Found}

\subsection{Critical errors}

\subsubsection{Synthetic experiment in code does not match the paper's stated simulator}

\status{Critical mismatch}

\paragraph{Paper statement.}
The paper says the synthetic experts are
\[
\hat y_{t,k} = 0.8 y_{t-1} + b_k + \varepsilon_{t,k},
\]
with regime-dependent correlated additive noises constructed from two shared
scalar noises and idiosyncratic noise; see
\code{Paper/icml2026/Section/Experiments.tex:67--90}. It also states
\mbox{$d_g = 1$} in \code{Paper/icml2026/Section/Experiments.tex:119--123}.

\paragraph{Code reality.}
The synthetic config instead sets
\code{shared_noise_mode: "model_match"} and \code{shared_dim: 2} in
\code{config/exp_synthetic_1.yaml:44--46} and
\code{config/exp_synthetic_1.yaml:78--79}. The environment then generates
latent shared factors $g_t$, latent expert-specific states $u_t$, and residuals
via
\[
\texttt{residuals[t] = x\_t * (loadings @ g\_t + u\_t) + noise},
\]
implemented in \code{environment/synthetic_env.py:571--577}. The actual expert
prediction is
\[
\hat y_{t,k} = y_t + \texttt{residual}_{t,k}
\]
in \code{environment/synthetic_env.py:841--846}.

\paragraph{Why this matters.}
This is not a cosmetic mismatch. The code's simulator is substantially closer to
the proposed model class than the paper's additive-intercept simulator. That can
materially bias the benchmark in favor of L2D-SLDS.

\paragraph{Minimal fix.}
Either:
\begin{itemize}
    \item rewrite the paper to describe the actual latent residual generator, or
    \item change the experiment code to the stated additive-noise simulator and
    rerun the results.
\end{itemize}

\subsubsection{Context-dependent transitions are random and untrained in synthetic and Melbourne runs}

\status{Critical paper-code mismatch}

\paragraph{Paper statement.}
The method is presented as using context-dependent regime transitions
$\Pi_\theta(x_t)$ with learned low-rank attention logits; see
\code{Paper/icml2026/Section/Approach.tex:135--163}.

\paragraph{Code fact.}
The \code{FactorizedSLDS} defaults are
\code{transition_mode="attention"} and \code{transition_init="random"} in
\code{models/factorized_slds.py:114--115}. If no trained transition model is
present, \code{\_context\_transition} constructs transition scores from those
parameters in \code{models/factorized_slds.py:741--763}. The synthetic config
uses \code{transition_mode: "attention"} and \code{em_enabled: false} in
\code{config/exp_synthetic_1.yaml:98--128}. The Melbourne config does the same
in \code{config/config_melbourne.yaml:123--153}.

\paragraph{Consequence.}
In those experiments, the purported context-aware transition kernel is not
learned. It is simply a random function of context unless some external
initialization path exists. I found no such training path for those runs.

\paragraph{Why this matters.}
This invalidates the natural reading that the reported no-EM results are using a
meaningful learned context-driven switching mechanism.

\paragraph{Minimal fix.}
If EM is disabled, default to a fixed homogeneous transition matrix or disclose
that the transition scores are random and not learned.

\subsubsection{FRED uses full-feedback offline EM and masked post-prefix reporting}

\status{Critical evaluation-protocol mismatch}

\paragraph{Code fact.}
The FRED config sets
\code{em_enabled: true}, \code{em_offline_full_feedback: true}, and
\code{em_tk: 2000} in \code{config/exp_FRED.yaml:166--168}. The driver forces
\code{router.feedback_mode = "full"} during offline EM in
\code{slds_imm_router.py:1983--2069}. The evaluation code masks all costs before
\code{em_tk} to \code{NaN} in \code{plot_utils.py:1085--1155}, and then reports
\code{np.nanmean} in \code{plot_utils.py:1291--1335}.

\paragraph{Consequence.}
The reported averages in such experiments are not full-horizon online
partial-feedback performance. They are post-warmup averages after a full-feedback
parameter estimation phase on the initial prefix.

\paragraph{Why this matters.}
This is a different experimental protocol. It may be perfectly legitimate, but
it must be stated plainly because it changes what the results mean.

\subsubsection{The paper's NeuralUCB baseline is not the implementation}

\status{Critical baseline mismatch}

\paragraph{Paper statement.}
The appendix describes NeuralUCB using gradient features
$h_{t,k} = \nabla_\omega f_\omega(x_t,k)$ and a single Gram matrix in parameter
space; see \code{Paper/icml2026/Section/AppendixParts/Experiments.tex:83--110}.

\paragraph{Code fact.}
The implementation explicitly states:
``This is a `NeuralLinear' approximation to NeuralUCB'' in
\code{models/neuralucb_baseline.py:22--25}. It uses:
\begin{itemize}
    \item a one-hidden-layer ReLU MLP embedding,
    \item per-expert ridge matrices in embedding space,
    \item SGD updates to the embedding.
\end{itemize}
See \code{models/neuralucb_baseline.py:70--83} and
\code{models/neuralucb_baseline.py:165--192}.

\paragraph{Consequence.}
The baseline is mislabeled. This matters because NeuralUCB's theoretical and
empirical behavior is not the same as a NeuralLinear approximation.

\subsection{Moderate concerns}

\subsubsection{Synthetic warmup described in paper is not implemented in the published config}

\status{Moderate mismatch}

\paragraph{Paper statement.}
The synthetic experiment section says: ``We simply run a small warmup of 100
steps before running L2D-SLDS and UCBs'' in
\code{Paper/icml2026/Section/Experiments.tex:121--123}.

\paragraph{Code fact.}
The synthetic config disables EM in \code{config/exp_synthetic_1.yaml:98}. The
main interaction loop starts from $t=1$ in \code{router_eval.py:514--554}. I did
not find a separate 100-step synthetic warmup path for the published config.

\paragraph{Interpretation.}
Either the text is outdated or the config is.

\subsubsection{Synthetic multi-seed reporting uses a fixed data seed by default}

\status{Moderate reporting issue}

\paragraph{Code fact.}
The synthetic environment hard-codes
\code{DATA_GENERATION_SEED = 42} unless \code{data_seed} is explicitly provided;
see \code{environment/synthetic_env.py:3--8} and
\code{environment/synthetic_env.py:44--50}. The helper script
\code{scripts/run_exp_synth_seeds.sh:34--39} only changes
\code{environment.seed}, not \code{data_seed}.

\paragraph{Consequence.}
By default, multi-seed synthetic runs vary learner randomness while keeping the
underlying synthetic time series and expert behaviors fixed.

\paragraph{Why this matters.}
That is acceptable if disclosed. It is misleading if the paper suggests the
error bars are over independent environment draws.

\subsubsection{Undocumented initialization heuristics affect Melbourne}

\status{Moderate method-reporting gap}

\paragraph{Code fact.}
The Melbourne config uses:
\begin{itemize}
    \item \code{init_R_mode: empirical_expert},
    \item \code{init_state_from_residuals: u},
\end{itemize}
in \code{config/config_melbourne.yaml:112--121}. These heuristics are implemented
in \code{slds_imm_router.py:1827--1951}.

\paragraph{Consequence.}
The model is not being evaluated only with the generic priors described in the
paper. It receives data-dependent initialization help.

\paragraph{Why this matters.}
Again, this may be sensible engineering, but it belongs in the paper if the
reported results depend on it.

\subsubsection{Diagnostic sample count can change the actual routing behavior}

\status{Moderate implementation quirk}

\paragraph{Code fact.}
When diagnostics are enabled, \code{select_expert} precomputes
\code{iz_scores} using \code{exploration_diag_samples} and then reuses those
scores in the action selection path; see
\code{models/factorized_slds.py:1820--1836}.

\paragraph{Consequence.}
The number of diagnostic Monte Carlo samples can affect the selected actions,
not just the logs. This is undocumented.

\subsection{Minor issues / style / clarity issues}

\begin{enumerate}
    \item \textbf{Wrong proposition cited in synthetic analysis.}
    The synthetic results section cites \code{prop:transfer} for query-time
    information transfer in
    \code{Paper/icml2026/Section/Experiments.tex:143--144}. The relevant result
    is \code{prop:cross_update}, not \code{prop:transfer}.

    \item \textbf{Melbourne text/table inconsistency.}
    Table values show NeuralUCB as \mbox{$5.88 \pm 0.04$} and Random as
    \mbox{$6.71 \pm 0.03$} in
    \code{Paper/icml2026/Section/Experiments.tex:177--182}, but the narrative in
    \code{Paper/icml2026/Section/Experiments.tex:194--196} says NeuralUCB is
    \mbox{$6.71 \pm 0.03$}.

    \item \textbf{Incorrect percentage arithmetic.}
    The increase from $5.69$ to $5.78$ is about $1.58\%$, not $16\%$.

    \item \textbf{Indexing inconsistency in synthetic regimes.}
    The text defines $z_t \in \{1,2\}$ in
    \code{Paper/icml2026/Section/Experiments.tex:56--57}, but the figure caption
    refers to ``Regime-0'' in
    \code{Paper/icml2026/Section/Experiments.tex:37--42}.

    \item \textbf{README command is stale.}
    The README points to \code{config/config_synth_paper.yaml} in
    \code{README.md:24}, but that file does not exist in the repository root.
\end{enumerate}

\section{Concrete Fixes}

\subsection{Paper fixes}

\begin{itemize}
    \item Rewrite the synthetic experiment description to match the actual code,
    or replace the code with the stated simulator and regenerate the table and
    figure.
    \item Correct the Melbourne table narrative, the 16\% arithmetic mistake, the
    regime indexing inconsistency, and the wrong proposition citation.
    \item State explicitly when experiments use offline EM, whether it uses full
    feedback, and whether reported averages exclude the warmup prefix.
    \item Replace ``NeuralUCB'' by ``NeuralLinear-style baseline'' unless the
    true algorithm is implemented.
    \item Add a short implementation-details paragraph documenting empirical
    initialization of $R$ and initial state covariances.
\end{itemize}

\subsection{Code fixes}

\begin{itemize}
    \item If \code{em_enabled: false}, set a deterministic fallback transition
    matrix instead of random attention weights, or require explicit transition
    parameters in the config.
    \item Add a paper-faithful synthetic mode implementing the exact equations in
    the manuscript.
    \item Decouple diagnostics from decision-making by ensuring
    \code{exploration_diag_samples} never overrides the action-path Monte Carlo
    budget.
    \item Update the README command to a valid config path.
\end{itemize}

\subsection{Experimental fixes}

\begin{itemize}
    \item Rerun synthetic experiments under the exact simulator stated in the
    paper.
    \item Report both full-horizon and post-warmup averages whenever EM warmup is
    used.
    \item For synthetic seed sweeps, vary both the learner seed and the data seed
    unless the goal is specifically to condition on one environment draw.
    \item Either implement the baseline described as NeuralUCB or relabel the
    existing results.
\end{itemize}

\subsection{Theoretical fixes}

\begin{itemize}
    \item Narrow the framing: present the method as a model-based routing
    heuristic under partial feedback, not as a learning-theoretic result unless a
    regret statement is added.
    \item Add a discussion of identifiability and approximation error from
    posterior projection.
    \item State explicitly the minimal conditions under which shared-factor
    covariance yields residual-level transfer.
\end{itemize}

\section{Final Verdict}

\subsection{Is the problem formulation clear?}

Mostly yes. The paper defines the sequential routing problem, the observables,
the latent variables, the cost, and the one-step benchmark clearly. The main gap
is that it does not define an online-learning benchmark such as regret.

\subsection{Is the method mathematically sound?}

The core state-space model and filter are mathematically coherent. The main
approximation, namely the projection to a factorized posterior, is transparent
and honestly described. The weakest mathematical points are the absence of
identifiability discussion and the lack of any bound on approximation error or
decision-quality loss from the projection.

\subsection{Are the theoretical claims justified?}

Only partially. The structural propositions are correct but limited. The paper
does not justify stronger claims about online optimality, regret, or asymptotic
performance.

\subsection{Does the code faithfully implement the paper?}

The core \code{FactorizedSLDS} implementation is substantially faithful to the
paper's filtering and MCEM description. The \emph{experimental protocols} are
not faithfully aligned. The major paper-code mismatches are:
\begin{itemize}
    \item the synthetic simulator,
    \item the untrained random transition model in no-EM runs,
    \item the FRED full-feedback offline EM plus masked evaluation,
    \item the mislabeled NeuralUCB baseline.
\end{itemize}

\subsection{Top three risks that could invalidate the results}

\begin{enumerate}
    \item \textbf{Synthetic benchmark mismatch.} The code appears more
    model-matched to L2D-SLDS than the paper discloses.
    \item \textbf{Random transitions in no-EM experiments.} The claimed
    context-aware switching mechanism is not meaningfully learned in synthetic
    and Melbourne under the published configs.
    \item \textbf{Protocol and baseline mismatches.} Full-feedback EM on FRED and
    the mislabeled NeuralUCB baseline weaken the credibility of comparative
    empirical claims.
\end{enumerate}

\paragraph{Bottom line.}
As a modeling paper, the work has a coherent core. As a rigorous paper-code
package, it currently has significant specification and evaluation problems that
must be corrected before the empirical claims can be treated as fully reliable.

\end{document}
